\documentclass[11pt]{article}
\usepackage[margin=1in]{geometry}
\usepackage{xcolor}
\usepackage{hyperref}
\frenchspacing
\hypersetup{
  pdftitle={Response to Reviewer Comments},
  pdfauthor={Haihua Mo and Zhengyou Liang},
  pdfsubject={Response to reviewer comments on the revised manuscript}
}
\setlength{\parindent}{0pt}
\setlength{\parskip}{0.55em}
\setlength{\emergencystretch}{1.5em}
\newcommand{\reviewcomment}[1]{\vspace{0.4em}\noindent\textbf{Reviewer comment.} \emph{#1}}
\newcommand{\response}{\vspace{0.25em}\noindent\textbf{Response.} }
\newcommand{\completed}{\vspace{0.25em}\noindent\textbf{Completed work and principal evidence.} }
\newcommand{\scope}{\vspace{0.25em}\noindent\textbf{Scope not evaluated and rationale.} }
\newcommand{\changes}{\vspace{0.25em}\noindent\textbf{Changes in the revised manuscript.} }

\begin{document}

\begin{center}
{\Large\bfseries Response to Reviewer Comments}\\[0.6em]
{\large Latency Optimization of Cascaded Voice Dialogue Systems with a Pipeline-Parallel Streaming Architecture}\\[0.4em]
Revision date: August 22, 2026
\end{center}

Dear Editor and Reviewer,

Thank you for the constructive assessment and for identifying the main evidentiary gaps in the original submission. We revised the manuscript to add human-recorded speech tests, perturbation and repeatability analyses, a configured streaming comparator, state-correction instrumentation, exploratory response-level evaluation, and a unified server-side first-playable-PCM measurement. We also narrowed claims that were stronger than the evidence.

The targeted experiments described below were completed during the revision. A subsequent final audit and deterministic CPU-only reanalysis corrected the interpretation and statistical treatment of the locked records. No new ASR, LLM, TTS, CUDA, audio-generation, corrected-trigger, controlled-contention, or injected-jitter experiments were run during this audit and reanalysis. We therefore distinguish explicitly between requests addressed experimentally and conditions that remain outside the present systems evaluation. All locations below refer to physical pages in the current 15-page compiled revised manuscript.

\section*{Comment 1: Human-recorded speech and robustness}

\reviewcomment{The main benchmark is synthesized from dialogue text, and the non-streaming ASR consequently has near-zero error. Please validate the system on real recordings with speaker, accent, noise, speaking-rate, and endpoint variation; these conditions directly affect VAD, timestamp stability, and transcript commitment.}

\response We addressed this concern in part by adding public human-recorded speech from LibriSpeech test-clean and AISHELL-1, together with controlled pause, noise, babble, and speaking-rate variation. These tests examine whether the latency behavior of the historical implementation persists outside the TTS-generated benchmark. The resulting test sets are concatenated human-read stress workloads, not spontaneous or natural dialogue recordings.

\completed On the second platform, each language contributes 75 clean samples (30 Long, 30 Very Long, and 15 Extra Long). Six additional 30-sample conditions per language use SNRs of 20, 15, and 10~dB, speed factors of 0.9 and 1.1, and 15-dB babble, with inserted pauses drawn from 0.2--1.0~s to provide controlled pause variation. For Extra Long clean speech, mean post-feed residual TTFT is reduced by 67.5\% on LibriSpeech and 65.7\% on AISHELL-1. We do not claim that recognition quality is preserved: mean utterance-level WER is 2.97\% for System~A and 5.65\% for System~B on LibriSpeech, while mean utterance-level CER is 10.77\% for System~A and 11.80\% for System~B on AISHELL-1. Babble is a failure boundary. For English babble, System~B is 47.3\% slower in the unpaired summary, but the paired analysis of 29 common samples yields a 95\% confidence interval that includes zero. For Chinese babble, the 6.5\% reduction is not significant ($p=0.271$). Empty streaming outputs occur in 12/30 English and 5/30 Chinese babble samples, and the slowest queue requires approximately 21~s after audio completion to drain. A five-sample manual spot check found all five samples intelligible, with no truncation or ordering errors, while two samples had audible concatenation seams.

\scope The added corpora broaden speaker and recording diversity, but we did not perform accent-stratified analysis, directly quantify timestamp drift, or repeat the commitment instrumentation on the human-read and noisy-speech sets. The inserted pauses provide controlled artificial endpoint variation rather than naturally occurring conversational turn endings. We also did not add spontaneous conversation, a study of microphone or channel effects, task-success testing, a noninferiority margin for response quality, or a user study. Those items require a separate data-collection and task-evaluation program. The revised manuscript therefore uses these data as a controlled non-TTS robustness check and does not claim robustness across accents, field conditions, or natural dialogue.

\changes The construction protocol is reported in Section~IV-D5, ``Concatenated Human-Read Speech and Augmentation'' (pp.~8--9). Results appear in Section~V-D and Table~VI (pp.~10--11). The historical trigger boundary and the remaining limitations on ecological validity are stated in Section~IV-B (pp.~5--6) and Sections~VI-B--C (p.~13).

\section*{Comment 2: Server-side time to first playable PCM}

\reviewcomment{TTFT is measured from the end of speech to the first LLM token and excludes network and TTS latency. Since the paper discusses voice-response latency, please also report time to first audible output, including endpointing, communication, and initial TTS generation.}

\response We added a unified server-side measurement from the estimated physical end of the prerecorded waveform to the first availability of playable PCM. The endpoint is named \emph{server-side time to first playable PCM}, rather than user-perceived audible latency, because no client playback or acoustic-onset timestamp was recorded. A single monotonic clock is used to timestamp pipeline closure, first-token generation, policy-defined response-text readiness, the localhost TTS request, and initial synthesis.

\completed On the second platform, the playable threshold is 1324 bytes, corresponding to approximately 30~ms of 22.05-kHz mono 16-bit PCM. The six timestamp components close exactly for all 100 mode--sample records (50 per mode). The two systems use different response-generation and TTS policies: System~B sends its first completed sentence, whereas System~A sends its capped full response. Under this policy bundle, median server-side time to first playable PCM is 3.11~s for System~B and 22.27~s for System~A. This is not a matched-text, matched-response-policy, or architecture-only comparison. Mean TTS text lengths are 204.10 characters for System~A and 17.74 characters for System~B, with exact greetings in 0 of 50 versus 25 of 50 cases, max-token stops in 41 of 50 versus 43 of 50 cases, and 19 of 25 English-input System~B TTS texts containing no ASCII Latin letters. The latter is reported only as a deterministic descriptive proxy, not as a language-quality metric. ASR and the localhost CosyVoice service share GPU0, while Qwen2 runs on GPU1. ASR--LLM communication uses in-process queues, and TTS requests are sent over localhost. The mean 133.0-ms feed-end-to-input-close component is a complete waiting interval, not flush-compute time, and explicit flush itself is approximately 0.21~ms.

\scope We did not measure external client transport, a playback jitter buffer, audio-device scheduling, DAC/speaker delay, or acoustic onset at the listener. The physical-speech-end anchor is estimated from the prerecorded waveform and mapped to the replay clock, so it is not a live conversational end-of-turn decision. We also did not run a distributed-network experiment, a whole-system matched-text comparison, or a policy-matched $2\times2$ design crossing serial versus pipeline scheduling with first-sentence versus full-response synthesis. These would require a distinct client-and-device harness and additional generation and TTS experiments. Rather than infer unmeasured quantities, the manuscript limits the result to the measured single-machine server-side policy comparison.

\changes Endpoint definitions appear in Section~I (pp.~1--2) and Section III-C1 (p.~4). Platform and policy details appear in Section V-A (p.~9). The unified measurement and Table~VIII are in Section V-G (pp.~12--13), and the client-side and policy-matching boundaries are reiterated in Sections VI-A--C (pp.~12--13).

\section*{Comment 3: Variability, percentiles, and repeated measurements}

\reviewcomment{Tables III--V mainly report mean values. Add standard deviations or percentile latency, especially P95/P99, and repeat measurements to show whether the approximately 1.1 s plateau remains stable under GPU contention and segment-trigger jitter.}

\response We added sample standard deviations, P95 and P99 tail latencies, observed P5--P95 bands, repeated measurements, and dialogue-cluster uncertainty estimates. These additions characterize variability in the recorded laboratory conditions. We did not perform stress tests with controlled external GPU contention or injected segment-trigger jitter, and we therefore do not claim stability under those conditions.

\completed Tables~III--V now report mean, sample standard deviation, P95, and P99 where applicable. Figure~6 shows 12 equal-frequency duration bins with observed P5--P95 bands, not confidence intervals. From Long to Extra Long, System~B's P99 grows to 1.32 times its Long-condition value, whereas System~A's grows to 4.96 times its Long-condition value over the measured duration range. This is weaker growth rather than a constant or universal bound. On the second platform, the repeat study uses 50 Very Long accumulated turns from 40 source dialogues and three runs per mode. System~B round means are 1434.0, 1482.1, and 1454.9~ms. The per-sample CVs have a median of 4.05\%, a P90 of 10.73\%, and a maximum of 18.96\%. Because accumulated turns share dialogue history, the locked primary inference uses 498 valid turns from 99 dialogues and 10,000 dialogue-cluster bootstrap resamples. On the original platform, System~B's ratio-of-means reduction relative to System~A on this set is 74.34\%, with a 95\% CI of [72.62\%, 75.76\%]. Cluster sizes range from 1 to 9 turns per dialogue, with a median of 5. The candidate ledger contains 505 turns. Run-log review identified seven executions contaminated by concurrent external programs, leaving the 498-turn valid cohort. Those seven records are retained in the audit ledger but are treated neither as system outcomes nor as evidence from a controlled contention experiment.

\scope All reported streaming measurements use the historical implementation's latched startup-duration policy: the cumulative admitted-duration counter is not decremented after processed segments are removed, so the 2.0-s condition is a one-time startup gate rather than a recurring current-queue-duration threshold. We did not evaluate a corrected trigger, inject trigger timing or arrival jitter, counterbalance the historical arm order, or run controlled contention experiments involving unrelated GPU jobs, shared services, or multiple users. No unrelated GPU jobs were present in the retained valid runs, and for R7 specifically, ASR and TTS were intentionally co-resident on GPU0. Corrected-trigger, counterbalanced, and controlled-stress comparisons would require new model runs. The revised manuscript therefore confines its conclusion to the measured fixed-order historical policy and lists the unmeasured conditions as future work.

\changes The statistical method is in Section III-C2 (p.~4); the historical trigger is documented in Section IV-B (pp.~5--6); execution order and platform conditions are in Section V-A (p.~9). Figure~6, Section V-B, and Table~III report the duration and repeatability results (pp.~9--10). The 505-to-498 ledger, cluster analysis, and fixed-order arm contrasts are in Section V-C and Table~IV (pp.~10--11). Remaining contention, jitter, order, and corrected-trigger limits are in Sections VI-B--C (p.~13).

\section*{Comment 4: Configured stronger streaming comparator}

\reviewcomment{The comparison is limited to the authors' non-streaming System A. Include a stronger streaming baseline or alternative chunked-ASR/prefill implementation, with matched models and hardware, so the gain can be attributed to the proposed scheduling and cache design.}

\response We added a configured, project-internal LocalAgreement-2 (LA-2)-style streaming comparator. It uses the same Whisper and Qwen2 model weights, inference engine, audio segmentation, second-platform hardware, and locked 498-turn/99-dialogue analysis set as System~B. The trigger policies are not matched: System~B uses the historical latched startup gate, whereas the configured LA-2-style implementation follows its own new-audio/agreement policy. The result therefore characterizes one configured operating point and does not establish general superiority over the LocalAgreement family or isolate an architecture-only causal effect.

\completed On the second platform, the comparator uses absolute-time-axis commitment, punctuation-robust agreement, sentence-boundary-aware trimming, and a 15-s retained-buffer cap. At this operating point, its mean post-feed residual TTFT is 2115.0~ms, compared with 1573.9~ms for System~B. The paired mean difference, calculated as the configured LA-2-style comparator minus System~B, is 541.11~ms, with a 95\% dialogue-cluster bootstrap CI of [485.66, 599.14]~ms and a Holm-adjusted dialogue-level $p=1.14\times10^{-17}$. Recognition errors differ across systems and languages, so no quality equivalence is claimed. On the original platform, the historical three-arm analysis further shows that most of the measured fixed-order contrast is associated with the streaming-ASR arm transition: Baseline minus Streaming-ASR-Only is 3332.13~ms, with a 95\% CI of [3064.38, 3572.61]~ms, whereas Streaming-ASR-Only minus Full-Streaming is 15.50~ms, with a 95\% CI of [10.95, 19.66]~ms. These are fixed-order arm contrasts, not causal component effects, and a stable large independent incremental-prefill benefit is not established under Qwen2-7B. Related Work was also expanded with direct references to Simul-Whisper, Prompt Cache, and SGLang with its RadixAttention mechanism.

\scope We did not sweep multiple LocalAgreement operating points, construct a latency--quality frontier, match the two trigger policies, evaluate a corrected-trigger System~B, add another streaming-ASR family benchmark, or rerun the historical three arms in counterbalanced order. Those additions require substantial new model experiments. A single stronger comparator, evaluated at one configured matched-model and matched-hardware operating point, addresses the request in part, and the manuscript responds to the remaining uncertainty by narrowing the claim rather than generalizing beyond the evidence.

\changes The expanded discussion of streaming ASR and cache-related research appears in Section~II (p.~2). The fixed-order arm contrasts appear in Section~V-C and Table~IV (pp.~10--11). The comparator configuration and results appear in Section~V-E (pp.~11--12), with Table~VII on p.~12. The bounded interpretation appears in Sections~VI-A--C (pp.~12--13), and the additional references are listed on pp.~14--15.

\section*{Comment 5: Commitment, tokenizer boundaries, KV state, and response quality}

\reviewcomment{Please clarify how committed text, tokenizer boundaries, and the KV cache are handled when a partial ASR hypothesis changes. Report rollback or correction frequency and evaluate downstream response quality, because WER/CER alone does not show whether incremental prefilling preserves dialogue meaning.}

\response We now distinguish internal ASR hypotheses from the downstream commitment contract. An internal hypothesis may change when retained audio is re-recognized, but the downstream delivered stream is append-only, and committed fragments are not withdrawn. The current prototype has no KV-cache rollback or repair path. Later corrections are logged internally but are not propagated, so the originally committed fragment remains in the downstream cache and may influence generation. We instrumented internal corrections and tokenizer seams and added a limited exploratory response-level analysis. These measurements do not establish semantic, task, KV-state, output-distribution, or usability equivalence.

\completed On 50 fixed Very Long accumulated turns from 40 source dialogues, instrumentation records 375 downstream commit-delivery events covering 425 committed segment IDs. Downstream rollback delivery is zero by construction because rollback is not implemented. There are 224 internal correction events, affecting 224 distinct committed segment IDs (52.7\%) and 49/50 turns. Correction edit distance has a mean of 2.3 characters, a P90 of 6, and a maximum of 16. Fragment-wise and one-shot token sequences differ in 25/50 turns, and both sequences decode to identical text in 50/50 cases. This establishes only text reconstruction, not equivalence in token sequences, KV states, logits, response distributions, or downstream effects. The exploratory response analysis uses independently sampled 128-token-capped replies. The mean BGE-M3 cosine similarity is 0.8832. In a single blinded intent-satisfaction evaluation, the service-hosted judge assigns mean scores of 3.10 out of 5 to System~A and 3.04 out of 5 to System~B, and the paired System~B-minus-System~A difference is $-0.06$, with a bootstrap 95\% CI of [$-0.34$, 0.22]. We report the service identifier submitted with the request and the access period, but the service exposed no immutable backend build identifier.

\scope We did not implement rollback-aware cache repair or retokenization and replay after correction. We also did not measure emission-time committed WER/CER, classify corrections involving entities, numbers, dates, negation, or intent, use deterministic paired response generation, evaluate task success, obtain human ratings, conduct repeated or multi-judge assessments, or prespecify a noninferiority margin for response quality. These would require further implementation work and a task-level evaluation program. The present revision directly reports how the implemented state contract behaves and quantifies its correction and seam boundaries, and the response-level analysis remains exploratory rather than evidence of meaning preservation.

\changes The general cache distinction is described in Section~III-B2 (p.~4), the commitment and seam contract in Section~IV-B4 (p.~7), and incremental cache behavior in Section~IV-C (pp.~7--8). The correction, seam, and exploratory response measurements appear in Section~V-F (p.~12), and their evidentiary limits are reiterated in Sections~VI-A--C (pp.~12--13).

\section*{Closing statement}

These revisions address the principal concerns and make explicit the requested conditions that remain outside the present evaluation: spontaneous and accent-stratified speech, live client-side endpointing and playback, matched response/TTS policies, controlled external contention and trigger jitter, corrected-trigger and counterbalanced reruns, a LocalAgreement operating-point frontier, rollback-aware cache repair, and larger-scale, repeated task-level quality assessments. We have limited the claims accordingly rather than inferring results for experiments that were not conducted.

Thank you again for the detailed review, which substantially improved the precision, transparency, and reproducibility of the manuscript.

Sincerely,

Haihua Mo and Zhengyou Liang

\end{document}
